\subsection{Addition as normalization: from normal forms to a commutative monoid}\label{subsec:emergent-arithmetic-stable-add}

The preceding definition of \(\oplus\) is algebraic transport along the bijection \(\mathrm{Val}:X_\infty\to\NN\). The point of the present subsection is to separate that tautological level from the genuinely digit-level statements about normalization and finite-state realizability.

\begin{remark}[A clarification about ``tautology'': definitional isomorphism versus observable implementation]\label{rem:arith-not-novel-as-abstract}
One should acknowledge that Definition \ref{def:stable-add} transports the operation directly from \((\NN,+)\). In that sense, the statement that \(\mathrm{Val}\) is an isomorphism is \textbf{definitional} at the level of abstract algebra: it does not produce a monoid that is nonisomorphic to \((\NN,+)\).

The reason for recording the isomorphism anyway is twofold. First, it fixes notation and emphasizes that addition need not be treated as a primitive input. Second, the real mathematical content lies in two further statements that interface directly with the ``scan and read out'' viewpoint:
\begin{itemize}
  \item \textbf{Orbit/folding content.} The congruence classes generated by local Fibonacci rewriting are exactly the value fibers, and the Zeckendorf language supplies the unique transversal; this is the content of Theorem \ref{thm:monoid-quotient-is-N} and Corollary \ref{cor:fold-as-section}.
  \item \textbf{Digit-level content.} The operation may be described as digitwise superposition followed by normalization, as made explicit in Corollary \ref{cor:add-as-fold}. This turns the abstract formula \(\mathrm{Val}^{-1}\) into a concrete, auditable normalization principle.
\end{itemize}
\end{remark}

\begin{theorem}[Digit-level realization of Zeckendorf addition with finite delay]\label{thm:zeckendorf-add-fst}
Let \(c,d\in X_\infty\) have finite support, and let
\[
m:=\max\bigl(\{k:\ c_k=1\text{ or }d_k=1\}\cup\{0\}\bigr)
\]
be the effective resolution. Choose any padding length \(L\ge m+2\), append zeros through position \(L\), and read the finite input synchronously from high Fibonacci weight to low Fibonacci weight:
\[
(c_L,d_L),(c_{L-1},d_{L-1}),\dots,(c_1,d_1)
\in(\{0,1\}^2)^L,
\]
where \(c_i=d_i=0\) for \(i>m\). There are a finite set of states \(Q\), an initial state \(q_0\), a transition/output rule and final-output map
\[
\tau:Q\times \{0,1\}^2\longrightarrow Q\times\{0,1\}^*,
\qquad
\omega:Q\longrightarrow \{0,1\}^*,
\]
and constants \(D,B\), independent of \(m,c,d\), such that after reading through input position \(i\) the transducer has irrevocably output all result digits in positions at least \(i-D\). After the final pair \((c_1,d_1)\) is read, the final-output map \(\omega\) emits the remaining low-order suffix of length at most \(B\) and the computation halts. Hence \(c\oplus d\) is computed in \(O(m)\) digit operations without constructing \(\mathrm{Val}(c)\) or \(\mathrm{Val}(d)\).
\end{theorem}

\begin{proof}
This is the golden-ratio instance of Frougny's online-addition theorem for numeration systems with Pisot base: for such systems, digitwise addition followed by normalization is realized by an online finite automaton with bounded delay \cite{Frougny1999OnlineAddition}. In the present notation the base is the Pisot unit \(\phi\), the admissible language is the golden-mean language of words with no factor \(11\), the input alphabet for addition is the synchronous pair alphabet \(\{0,1\}^2\), and the preliminary digitwise sum lies in the finite alphabet \(\{0,1,2\}\). The cited online interface reads from most significant digit to least significant digit and supplies a delay bounded by a constant depending only on the numeration system, not on the input length.

Composing the letter map \((a,b)\mapsto a+b\) from \(\{0,1\}^2\) to \(\{0,1,2\}\) with this finite-state normalizer gives the required transducer \((Q,q_0,\tau,\omega,D,B)\). The normalizer is value preserving and returns the unique admissible Zeckendorf representative; by Corollary \ref{cor:add-as-fold}, that representative is exactly \(c\oplus d\). Leading zero padding does not change value, and the final-output map emits only a bounded suffix because the online delay is bounded and the numeration system is fixed. Since one transition is performed per input pair and only bounded final output remains, the running time is linear in \(m\). Concrete Zeckendorf implementations of the same finite-state normalization principle are given by Ahlbach--Usatine--Frougny--Pippenger \cite{AhlbachUsatineFrougnyPippenger2013}; the Ostrowski formulation is also recorded in \cite{HieronymiTerry2018OstrowskiAddition}.
\end{proof}

\begin{remark}[Low-to-high scans versus online numeration convention]\label{rem:online-orientation}
The stable-address indices in the rest of the paper are written from low weight to high weight because this is convenient for the value formula \(\sum c_kF_{k+1}\). The online theorem above uses the standard numeration-theoretic convention of reading most significant digits first. This is only an interface convention: for a finite support input one first chooses the effective resolution \(L\), scans downward from \(L\), and then interprets the output back in the same stable-address coordinates.
\end{remark}

\begin{remark}[Online realizability and non-parallelizability]\label{rem:online-not-parallel}
In the scan-based semantics, the existence of an online transducer means that readout may proceed unidirectionally in time while maintaining only finite memory. This fits the folding perspective adopted in the present paper. At the same time, addition in the golden-ratio numeration system, although online, is in general not parallelizable in the strong circuit-theoretic sense. In the present framework, this reflects the fact that stable readout has finite memory but still carries cross-scale dependence through carries and normalization. See \cite{Frougny1999OnlineAddition} for precise results and discussion.
\end{remark}
